\chapter{Introduction}
\section{Project Context}
Games have long served as benchmarks for measuring Artificial Intelligence (AI) progress \cite{StudeGame}. Many real-world problems, from military tactics to multiplayer games, require intelligent decision-making in environments where all information is not available. A common example is the ``fog of war'' mechanic found in many multiplayer games \cite{FogofWar}, where the game environment is not fully revealed to any single player, placing these games in the category of imperfect information games \cite{DRLImperfectInfo2023}. In such games, players must develop strategies without full knowledge of the game's environmental state \cite{FogofWar}. Determining the optimal strategy under these conditions remains a standing challenge for Artificial Intelligence \cite{li2024opponent}.

This study focuses on mastering imperfect information games using Multi-Agent Reinforcement Learning (MARL) with Deep Q-Networks (DQN). In MARL, multiple agents learn and adapt in the same environment, making each decision based on its own observations and rewards \cite{2509_03682v1}. With DQN, the agents use deep neural networks to approximate the Q-function that maps state-action pairs to expected reward. This study introduces the MARQ (Multi-Agent Rainbow Deep Q-Networks) framework, an extension of the DQN architecture with Rainbow enhancements: distributional value estimation (C51) and noisy exploration networks \cite{RainbowJMLR2021}. An AAREN (Attention as Recurrent Neural Network) module \cite{Aaren} processes extended observation sequences and produces learned embeddings meant to carry what the opponent's moves reveal about its hidden pieces. AAREN is trained on the pieces that battles reveal, and the Q-network reads its embeddings as inputs. League-based self-play training \cite{perolat2022stratego} introduces agents to different historical opponents to prevent similar strategies. Stratego is a good example of a test environment for our framework, agents compete against each other and must develop strategies through repetition of experience while reasoning under uncertainty about the enemy's setup and the environments hidden states.



\section{Purpose and Description}
This research evaluates whether Multi-Agent Reinforcement Learning (MARL) with Deep Q-Networks (DQN) can produce agents that can play imperfect information games at a competitive level. The board game Stratego serves as the test domain because it combines hidden piece identities, a game tree of about $10^{535}$ nodes \cite{perolat2022stratego}, and long game horizons. Existing analytical tools for Stratego have little to no support for the deeper tactical reasoning that we want to achieve. 

The MARQ framework was developed to test whether targeted architectural changes, specifically distributional value estimation, attention-based history encoding, and learned exploration, can overcome the performance ceiling observed in a standard DQN baseline. The experiments, detailed in later chapters, do not show that they did. In self-play the combined changes cut the draw rate from 50.04\% to 6.69\%, and the self-play win rate rose from 24.97\% to 49.88\%. Both sides of a self-play game are the same agent, however, so that rise follows from the fall in draws and does not measure strength (Section~\ref{finding:7}). Against fixed heuristic opponents the Rainbow agent with AAREN won 11.66\% of 75,047 games and drew 77.02\%. The runs made after those experiments, with AAREN switched off, found that the agent played at the level of the scripted player it had imitated before training, and that reinforcement learning did not raise that level (Appendix~\ref{app:lessons}). The thesis therefore reports what was built, what was measured, and a diagnosis of why the learner stalled.
\subsection{Research Questions}
This study is guided by three main research questions:

\begin{enumerate}
   
    \item How does the MARQ framework perform in terms of convergence speed and win-rate against heuristic baselines in a partially observable environment?
    \item Does the MARQ framework demonstrate statistically significant superiority over standard DQN and rule-based agents across varying opponent difficulty tiers?
    \item What is the impact of integrating the AAREN module on the agent's ability to infer hidden opponent piece configurations compared to approaches without explicit belief state modeling?
\end{enumerate}

Chapter~\ref{chap:recommendations} answers each question as far as the data allow.

\section{Objectives}
The main objective of this study is to develop and evaluate the MARQ framework for Stratego, where agents must operate under hidden state variables. The specific objectives are:
\begin{itemize}
    \item Identify and evaluate reinforcement learning algorithms for competitive play in Stratego.
    \item Design and implement a game environment for Stratego
    \item Build and implement a training model using Multi-Agent Reinforcement Learning with Deep Q-Networks and AAREN.
    \item Train and optimize the AI agent to respond to varying game states, hidden information, and opposing strategies.
\end{itemize}

\section{Scope and Limitation}

The study is limited to the development of an AI agent for imperfect information games, with Stratego as the test environment. The scope includes the design and implementation of Stratego as a game environment, the development of a training framework using Multi-Agent Reinforcement Learning (MARL) with Deep Q-Networks (DQN) and AAREN, and the training and optimization of the AI agent to adapt to varying game states. Agent performance is evaluated through agent-vs-agent testing within the Stratego environment. Each configuration was trained once, and no ablation that removes one component at a time was run, so the results describe the combined system and not the contribution of each part. Training is limited to the personal devices of the researchers and does not rely on cloud-based computing resources.

\section{Significance of the Study}
Current imperfect information game systems either require large computational resources or use model-free approaches with limited strategic depth. The MARQ framework contributes a specific combination of techniques: implicit opponent modeling through AAREN-generated embeddings, distributional value estimation via C51, and learned exploration via Noisy Networks. The experiments do not establish that this combination plays better than the standard DQN baseline: the two were never evaluated against each other or against the same fixed opponents, and against heuristic opponents the combined agent won 11.66\% of its games.

The study contributes three things that hold whatever the outcome. The first is a tested Stratego environment and training system on which a value-based learner can be run and examined on a single consumer GPU, with a network that can be played in a browser. The second is a set of checks that reversed results the study had first read as validation: a self-play win rate that measured the draw rate, a loss--win correlation that measured a shared trend, an average Q-value taken over illegal actions, an accuracy compared with the wrong baseline, and a gradient path that reached a single bias (Chapter~\ref{chap:results}, Appendix~\ref{app:lessons}). The third is a measured account of where the agent's strength came from: the demonstrator it imitated set its level, a behaviour cloning loss over the legal moves transferred the demonstrator's decisions where a large-margin loss had not, and the policy was lost when the cloning term was withdrawn (Appendix~\ref{app:lessons}, Section~\ref{sec:formal14}). These apply to other partially observable or adversarial domains in which a value-based learner is started from demonstrations.

\section{Definition of Terms}

\noindent\textbf{AAREN (Attention as Recurrent Neural Network)} \\
A neural architecture that combines the parallel training efficiency of Transformers with the sequential processing of RNNs. In the MARQ framework, AAREN processes observation sequences to produce learned embeddings that encode implicit piece identity inference. \\[6pt]

\noindent\textbf{Counterfactual Regret Minimization (CFR)} \\
An iterative algorithm that approximates Nash equilibrium by minimizing regret, widely applied in imperfect information games such as Poker and Stratego. \\[6pt]

\noindent\textbf{Deep Q-Networks (DQN)} \\
A reinforcement learning algorithm that combines Q-learning with deep neural networks, enabling agents to approximate optimal decision-making in complex environments. \\[6pt]

\noindent\textbf{Distributional Reinforcement Learning} \\
An approach that models the full distribution of returns rather than just the expected value. The MARQ framework uses C51 (Categorical DQN) with 51 atoms to capture uncertainty in value estimates. \\[6pt]

\noindent\textbf{Fog of War} \\
A game mechanic used in strategy games that conceals parts of the map or opponent actions, requiring players to utilize prediction and information-gathering strategies to reduce uncertainty. \\[6pt]

\noindent\textbf{Imperfect Information Games} \\
Games where players have incomplete knowledge of the current game state, requiring them to make strategic decisions under uncertainty (e.g., Stratego, Poker, Multiplayer Online Battle Arenas). \\[6pt]

\noindent\textbf{Information Set Monte Carlo Tree Search (IS-MCTS)} \\
A search algorithm for imperfect information games that samples consistent game states from belief distributions and performs tree search on information sets rather than individual states. \\[6pt]

\noindent\textbf{League Training} \\
A self-play training paradigm where agents compete against a diverse pool of historical opponents, preventing strategy cycling and ensuring robust generalization across different playing styles. \\[6pt]

\noindent\textbf{Multi-Agent Reinforcement Learning (MARL)} \\
An extension of reinforcement learning in which multiple agents interact within the same environment, learning coordinated or competitive strategies through simultaneous training. \\[6pt]

\noindent\textbf{Nash Equilibrium} \\
A game-theoretic concept where no player can gain an advantage by unilaterally changing their strategy, often used as a benchmark for decision-making in imperfect information games. \\[6pt]

\noindent\textbf{Perfect Information Games} \\
Games in which all players have complete knowledge of the current game state and available moves, such as Chess and Go. \\[6pt]

\noindent\textbf{AAREN Embeddings} \\
Embedding vectors produced by the AAREN module from observed opponent action sequences. \\[6pt]

\noindent\textbf{Rainbow DQN} \\
An enhanced DQN architecture that combines multiple improvements including distributional RL, noisy networks for exploration, and dueling network architecture for improved stability and sample efficiency. \\[6pt]

\noindent\textbf{Recurrent Neural Network (RNN)} \\
A type of neural network designed to process sequential data by maintaining memory of previous inputs, useful in imperfect information games for recalling past moves or revealed information. \\[6pt]

\noindent\textbf{Reinforcement Learning (RL)} \\
A machine learning paradigm where agents learn optimal strategies through trial-and-error interactions with the environment, guided by rewards and penalties. \\[6pt]

\noindent\textbf{Stratego} \\
A two-player strategy board game characterized by hidden information and asymmetric piece abilities, where the objective is to capture the opponent's flag or render them immobile. \\[6pt]